%!TEX root = ../Pflichtenheft.tex

% Kapitel 10
% Die Unterkapitel können auch in separaten Dateien stehen,
% die dann mit dem \include-Befehl eingebunden werden.
%-------------------------------------------------------------------------------

\chapter{Technical Product Environment}
\label{chap:tech_env}

This chapter describes the technical environment in which the product operates. It outlines the software architecture, technologies, and frameworks used for the implementation, as well as the tools and platforms that support development and deployment.

\section{Software}
The system is implemented in a \textbf{\hyperref[def:server-client]{Server-Client}} architecture, where the client application interacts with a centralized server to exchange data and perform operations.  
\subsection{Client}
The client component is a responsive single-page web-app that allows users to chat with an \hyperref[abr:AI]{AI} assistant, capable of generating code for a selection of \textit{\hyperref[abr:DSL]{DSL}s}, notably \textbf{\hyperref[abr:SPICE]{SPICE}}, based on a user prompt, and visualize generated code (if available). \newline
The client component will be developed using the  \textit{Angular Framework} \footnote{https://angular.dev/} with the \textit{Typescript} \footnote{https://www.typescriptlang.org/} programming language.

\subsection{Server}
The server component is a server that implements a \textit{\hyperref[abr:RAG]{RAG}} system and is programmed using the \textit{Python} \footnote{https://www.python.org/} programming language. \newline
The server component is connected to a \textit{Chroma} vector database for text embeddings, and a \textit{relational database} for storing app-specific data. The server component will use \textit{LangChain} and the \textit{Open\hyperref[abr:AI]{AI} \hyperref[abr:API]{API}} for communicating with the \textit{\hyperref[abr:LLM]{LLM}}.


\section{External interfaces}

\subsection{Graphical user interface}
The system will provide a \textit{Graphical user interface} (\hyperref[abr:GUI]{GUI}) as describes in~\ref{ch:ui}.
This \textit{\hyperref[abr:GUI]{GUI}} will be part of the client component and accessible via a browser. The interface enables users to: enter queries or prompts, view system responses or results, manage user-related actions such as login or session handling. The GUI enables users to enter prompts, view responses generated by the system, and manage application-specific features such as authentication or session control.

\subsection{Vector database}
The server component will interact with a \textit{\hyperref[def:vector-database]{Vector Database}} that will be used to store documents or specific \textit{\hyperref[abr:DSL]{DSL}} examples. Vector databases are used in \textbf{\hyperref[abr:RAG]{RAG}} systems to store vectorized text and query large amounts of documents based on similarity to a text snippet. Within this system, textual data is first transformed into high-dimensional vector embeddings, which are then indexed within the database. When a user submits a query, the system compares its embedding to stored vectors to identify the most semantically relevant information. This approach allows efficient handling of large text corpora and enhances the quality of responses provided by the language model.

\subsection{Large Language Model}
The server component will interact with a \textit{\hyperref[abr:LLM]{LLM}} that will provide the answers to the prompt provided by the user. This interaction enables the system to provide informative and context-aware answers that extend beyond simple pattern recognition.

\subsection{\hyperref[def:relational-database]{Relational Database}}
The server component will interact with a \textit{\hyperref[def:relational-database]{Relational Database}} (\hyperref[abr:RDB]{RDB}) to store application-specific data like user accounts, authentication tokens or conversations. 